\documentclass{article}
\usepackage{icmlko}
\usepackage{graphicx}
\usepackage{booktabs}
\usepackage{amsmath}

\title{같은 정답 위에서 다시 재다 \\
  {\large 19.9 분이 137.6 분이 되었고, 경계를 못으로 박고도 이겼다 ---
   그런데 천장은 여전히 세 배 위에 있었다}}

\begin{document}
\maketitle

\begin{abstract}
직전 노트는 순열 200 뽑기로 $p=0.004975$ 를 얻고 「순서는 잡힌다」를 확정했다.
프레시 티처가 그 판정의 \emph{설명}을 전부 무너뜨렸다 --- 귀무는 「꼬리가 짧은」 분포가 아니라
\textbf{네 점 혼합}이고 원인은 \texttt{prevmed} 3 분위 경계 하나이며(C1),
진짜는 200 뽑기가 한 번도 안 가는 \textbf{다섯째 상태}에 있고(C2),
무엇보다 판정에 쓴 판은 \textbf{진짜와 귀무가 서로 다른 정답 위에서 채점된다}(C4).
이 노트는 티처가 지목한 \textbf{깨끗한 판}에서 다시 쟀다.
\claimx{\textbf{정답 다중집합이 같은 판에서도 이겼다} --- [달력제거] 200 뽑기에서
$k=0$, $p=\mathbf{0.004975}$. 효과는 $\mathbf{137.6}$ 분이다.}
{직전 노트의 19.9 분은 \texttt{drop\_first} 가 \textbf{진짜 쪽에서만 긴 첫 간격}을 떼어 낸 결과였다
(진짜 중앙값 20 일 대 귀무 11 일). 정답을 맞추자 효과가 \textbf{6.9 배}로 드러났다.
정답 다중집합 동일성은 뽑기마다 찍었다 --- \textbf{200/200}.}
\claimx{🔴 \textbf{그리고 분위 경계를 못으로 박고도 이겼다} ---
진짜의 경계 $(4.5, 7.0)$ 을 200 귀무 \emph{전부}에 강제로 물려 이산 상태를 지웠는데
$k=0$, $p=\mathbf{0.004975}$, 효과 $\mathbf{131.1}$ 분.}
{티처의 「928 의 헤드라인은 분위 경계 하나의 이름 아니냐」에 대한 답은 \textbf{아니다}.
경계가 나르는 몫은 $137.6-131.1=\mathbf{6.5}$ 분뿐이다.
그리고 경계를 하나로 고정한 판에서는 상태 수가 1 이므로
\textbf{$z=13.80$ 이 이 계열에서 처음으로 해석 가능해진다}.}
\claimx{🔴🔴 \textbf{그런데 크기는 측정 전에 졌다} ---
채택 크기 $X=1.0$ 일 $\div$ 오라클 천장 $Y=0.3284$ 일 $=$ $Z=\mathbf{3.045}>1$.}
{$Z>1$ 은 \textbf{이 자로는 원리상 못 넘는다}는 뜻이다(규율 4).
그리고 진짜 팔은 이미 그 천장의 \textbf{98.60\%} 다 ---
\textbf{자료를 더 줘도 이 조건 열로는 못 올라간다.}
🔴 여는 방향은 하나뿐이다: \textbf{격자 정체}(이 판 오라클 $1.94386$ 일).}
\end{abstract}

\section{티처가 지목한 것을 그대로 했다}

직전 노트가 판정에 쓴 [둘 다] 판은 \texttt{drop\_first=True} 로 각 세계의 \emph{첫} 간격을 뗀다.
그런데 「첫 번째」는 세계마다 다른 행이다 --- 진짜에서는 진짜 첫 간격, 귀무에서는 무작위 간격.
러너의 \texttt{assert} 는 \textbf{격자 열만} 보증하므로 이것을 원리상 못 잡는다.

🔴 \textbf{깨끗한 판은 이미 저장소에 있었고, 판정에 안 썼다.}
이 노트는 그 판([달력제거] $\cdot$ \texttt{drop\_first=False})에서 다시 쟀다.

\begin{figure}[h]\centering
\includegraphics[width=\linewidth]{figs/fig1.pdf}
\caption{(a) 판 ① --- 경계 상태 4 개. (b) 판 ② --- 경계를 강제로 하나로 고정.
둘 다 $k=0/200$. (c) 귀무의 경계 상태별 개선.
진짜의 경계 $(4.5,7.0)$ 은 200 뽑기가 \textbf{한 번도 안 가는} 자리다.}
\end{figure}

\begin{table}[h]\centering\small
\begin{tabular}{lrr}
\toprule
& 판 ① [달력제거] & 판 ② [$\cdot$ 경계강제] \\
\midrule
진짜 & $0.3238034$ & 같음 \\
귀무 200 평균 & $0.2282464$ & $0.2327888$ \\
귀무 SD (ddof=1) & $0.0091586$ & $0.0065976$ \\
귀무 최대 & $0.2528891$ & $0.2543108$ \\
$k$ $\cdot$ $p$ & $\mathbf{0}$ $\cdot$ $\mathbf{0.004975}$ & $\mathbf{0}$ $\cdot$ $\mathbf{0.004975}$ \\
효과 & $0.09556$ 일 $=$ \textbf{137.6 분} & $0.09101$ 일 $=$ \textbf{131.1 분} \\
경계 상태 수 & 4 (\emph{$z$ 해석 불가}) & \textbf{1} ($z=13.80$) \\
\bottomrule
\end{tabular}
\caption{판 둘. 🔴 판 ① 의 $z\cdot$SD$\cdot$왜도$\cdot$첨도는 \textbf{혼합의 요약이라 해석하면 안 된다} ---
산출물이 그 사실을 필드로 들고 있다. 판 ② 에서 비로소 해석 가능해진다.}
\end{table}

\section{뽑기마다 관문을 찍었다}

직전 사이클들이 「호출했다」고 적고 안 부른 관문이 있었다(비교가능성 관문은 925$\cdot$928 에서
\textbf{호출 0 건}). 이번엔 \textbf{뽑기마다} 찍었다.

\begin{figure}[h]\centering
\includegraphics[width=\linewidth]{figs/fig2.pdf}
\caption{(a) 규율 4 --- $X/Y=3.045>1$ 이 \emph{측정 전에} 났다.
(b) 진짜 팔은 이미 천장의 98.60\%.
(c) 뽑기마다 찍은 관문 넷이 전부 200/200.}
\end{figure}

정답 다중집합 동일 \textbf{200/200} $\cdot$ 비교가능성 통과 \textbf{200/200}(상대차가 전부 정확히 $0.0$) $\cdot$
\texttt{mag} 경계가 진짜와 같음 \textbf{200/200} $\cdot$ 전체 중앙값이 진짜와 같음 \textbf{200/200}.
그리고 순열은 실제로 순서를 부순다 --- 간격 순서가 바뀐 격자 비율 중앙값 \textbf{86.1\%}.

\subsection{배선 --- 심어서 확인했다}

W5 는 티처 C4 가 지목한 자리에 \textbf{검정력이 있음을 증명한다}:
같은 두 세계를 \texttt{drop\_first=True} 로 보면 다중집합 검사가 \textbf{False} 를 내고,
[달력제거] 로 보면 \textbf{True} 를 낸다. W4 는 경계 강제에 $(\mathrm{nan},\mathrm{nan})$ 을 물려
개선이 $0.2178 \to 0.0468$ 로 무너지는 것을 보였다 --- \textbf{강제가 정말 먹힌다.}
회계: 심은 \textbf{6/6 발화} $\cdot$ 음성 대조 3/3 $\cdot$ 🔴 \textbf{못 심은 것 4}
$\to$ $6/6=1.000$ 과 $\mathbf{6/10=0.600}$ 을 \emph{둘 다} 신고한다.

\section{🔴 그런데 크기는 측정 전에 졌다 --- 그리고 그것이 이 노트의 결론이다}

규율 4 는 사전등록에 \textbf{「채택 크기 $X$ $\cdot$ 오라클 천장 $Y$ $\cdot$ $X/Y=Z$」}를 박게 한다.
$Z>1$ 이면 \textbf{그 사이클의 일은 측정이 아니라 자를 바꾸는 것}이다.

\begin{table}[h]\centering\small
\begin{tabular}{lrl}
\toprule
층 & 오라클 천장(개선) & \\
\midrule
기후값 상수 & $0$ & 🔴 $6.0$ 은 홀드아웃 최적 상수이기도 하다 \\
현 조건 열(dow$\cdot$quarter$\cdot$mag$\cdot$prevmed) & $\mathbf{0.32839}$ 일 & 진짜 팔이 여기의 \textbf{98.60\%} \\
\texttt{prevmed} 실수값 & $0.76978$ 일 & 아직 1.0 미만 \\
🔴 \textbf{격자 정체} & $\mathbf{1.94386}$ 일 & \textbf{1.0 을 넘는 유일한 층} \\
\bottomrule
\end{tabular}
\caption{홀드아웃 답을 알고 최적 적합한 오라클. 🔴 채택 크기 1.0 일은
현 조건 열의 천장보다 \textbf{3.045 배} 위에 있다.}
\end{table}

그러므로 이 노트가 낼 수 있는 정직한 문장은 둘이다.
\textbf{「순서는 잡힌다 --- 깨끗한 판에서, 경계를 고정하고도」}가 하나이고,
\textbf{「그러나 이 조건 열로는 채택 크기를 원리상 못 넘는다」}가 둘이다.
🔴 \textbf{그리고 둘째 문장은 문턱을 탓하는 것이 아니라 자를 갈라는 지시다.}

\section{🔴 자기 적발}

\begin{enumerate}
\item \textbf{사전등록 \S3 의 부등호가 틀렸다.} $2/201=0.00995<0.01$ 이므로
$p<0.01$ 의 동치 조건은 $k=0$ 이 아니라 $\mathbf{k\le 1}$ 이다.
러너가 그 대조를 필드로 갖고 있어 \textbf{정본에 \texttt{false} 가 박혔다} --- 사후 변명이 아니다.
판정은 불변($k=0$)이고, 신고한 검출력은 \textbf{하한}이었다(참값 $0.73759$ --- 정확히 두 배).
\item \textbf{연습 주행이 판정 씨앗 1$\sim$4 를 이미 돌았다.}
산출물을 \emph{안 지웠고} 본 주행과 \textbf{비트 동일 4/4} 를 대조했다.
\end{enumerate}

\section{진단 --- 판정이 아니다}

전체 차의 \textbf{86.3\%} 가 \textbf{첫 간격 행 2{,}277 개(4.15\%)}에서 온다.
거기서는 차가 $\mathbf{1.9424}$ 일로 $X=1.0$ 을 넘는다.
🔴 \textbf{그러나 축을 결과를 보고 갈랐으므로 결론에 안 넣는다}(사후 부분집합).
딱지도 함께 적는다 --- 그 부분집합에서는 격자 정체 오라클이 MAE $0.0$ 으로 \textbf{퇴화}한다
(격자당 첫 간격이 하나뿐이라서).

\end{document}
